% ==================================================================
\section{Chapter 5 --- Stochastic Integration}
% ==================================================================
Throughout this chapter, we argue on a filtered probability space $(\Omega, \mc{F}, (\mc{F}), P)$, and we assume that the filtration $(\mc{F}_t)_{t\geq0}$ is \textbf{complete}. Unless otherwise specified, all processes in this chapter are indexed by $\R_+$ and take real values. We often say "continous martingale" instead of "martingale with continuous sample paths". 

\subsection{Construction of stochastic integrals}
\subsubsection{Martingales bounded in $L^2$}
\begin{definition}[Space of integrators $\mb{H}^2$]\label{chap5.2:definition-space_of_integrators_H2}
\LG{Page 98} \\
We define $\mb{H}^2$ to be the space of all continuous martingales $M$ which are bounded in $L^2$ and such that $M_0$. We apply the usual convention that two indistinguishable processes are equal.
\end{definition}

\begin{proposition}[Characteristics of $\mc{H}^2$]\label{chap5.2:proposition-characteristics_of_integrators_H2}
\LG{Page 98} \\
$M \in \mb{H}^2$ if and only if $M$ is a continuous local martingale such that $M_0=0$ and $E[\br{M,M}_\infty]< \infty$. \\
Likewise if $M \in \mb{H}^2$, we have $M_t=E[M_\infty|\mc{F}_t]$ where $M_\infty \in L^2$ is the almost sure limit of $M_t$ as $t \to\infty$.
\end{proposition}
\begin{proof}
Use Proposition \ref{chap4.3:theorem-L2_bounded_martingales_and_quadratic_variation} and  Theorem \ref{chap3.4:theorem-martingale_equivalence_of_L1_convergence}.
\end{proof}


\begin{proposition}[Inner/scalar product and norm on $\mb{H}^2$] \label{chap5.2:proposition-H2_scalar_product_and_norm}
\LG{Page 98} \\
If $M,N \in \mb{H}2$, the random variable $\br{M,N}_\infty$ is well defined, and we have $E[|\br{M,N}_\infty|]< \infty$. The mapping $\mb{H}^2\times\mb{H}^2 \to \R: (M,N) \to (M,N)_{\mb{H}^2}$ given by
$$(M,N)_{\mb{H}^2}=E[\br{M,N}_\infty]=E[M_\infty N_\infty]$$
is a symmetric, bilinear form, where $(M,N)_{\mb{H}^2}=0$ if and only if $M=0$. The scalar product $(M,N)_{\mb{H}^2}$ thus yields a norm on $\mb{H}^2$ given by
$$||M||_{\mb{H}^2}=(M,N)_{\mb{H}^2}^{1/2}=E[\br{M,M}_\infty]^{1/2}=E[(M_\infty)^2]^{1/2}$$
\end{proposition}

\begin{proposition}\label{chap5.2:proposition-H2_Hilbert_Space}
\LG{Proposition 5.1}\\
The space $\mb{H}^2$ equipped with the scalar product $(M,N)_{\mb{H}^2}$ is a Hilbert space.
\end{proposition}


\begin{definition} 
An \textit{elementary process} is a progressive process of the form
$$H_s(\omega)=\sum_{i=0}^{p-1}H_{(i)}(\omega)1_{(t_{i}, t_{i+1}]}(s),$$
where $0=t_0<t_1<...<t_p$ and for every $i \in \{0,1,...,p-1\}$, $H_{(i)}$ is a bounded $\mc{F}_{t_i}$-measurable random variable. \\
Let $\mc{E}$ be the set of all elementary processes. 
\end{definition}


\begin{theorem}[Itô isometry / characterisation] \todo{$H\mapsto H\cdot M$ isometry $\LtwoM\to\HH$; $H\cdot M$ unique s.t.\ $\br{H\cdot M,N}=H\cdot\br{M,N}\ \forall N\in\HH$} \end{theorem}
\begin{proposition} \todo{associativity $K\cdot(H\cdot M)=(KH)\cdot M$; stopping $(H\cdot M)^T=H\ind_{[0,T]}\cdot M=H\cdot M^T$} \end{proposition}
\subsubsection{Local martingales}
\begin{theorem} \todo{extension to $\Lloc$ by localisation; same characterisation} \end{theorem}
\begin{remark} \todo{when is $H\cdot M$ a true martingale: $\E\int_0^tH_s^2d\br{M}_s<\infty$} \end{remark}
\subsubsection{Semimartingales}
\begin{definition} \todo{locally bounded progressive $H$; $H\cdot X=H\cdot M+H\cdot A$} \end{definition}
\subsubsection{Convergence of stochastic integrals}
\begin{proposition}[Dominated convergence] \todo{} \end{proposition}
\begin{proposition}[Riemann sums] \todo{$H$ adapted continuous: $\sum H_{t_{i-1}}(X_{t_i}-X_{t_{i-1}})\to\int H\,dX$ in probability} \end{proposition}
\begin{quizbox} The characterisation via $\br{H\cdot M,N}$ and the isometry — most proofs in 5.2--5.4 run through them. \end{quizbox}

\subsection{Itô's formula}
\begin{proposition}[Integration by parts] \todo{$X_tY_t=X_0Y_0+\int X\,dY+\int Y\,dX+\br{X,Y}_t$} \end{proposition}
\begin{theorem}[Itô's formula] \todo{$F\in C^2(\R^p)$, multidimensional} \end{theorem}
\begin{proofidea} \todo{IBP $\Rightarrow$ polynomials by induction $\Rightarrow$ Weierstrass + localisation; or Taylor + Riemann sums} \end{proofidea}
\begin{proposition}[Stochastic exponential] \todo{$\Exp(M)_t=\exp(M_t-\tfrac12\br{M}_t)$ is a nonneg.\ loc.\ mart / supermartingale} \end{proposition}
\begin{quizbox} Itô for $f(t,B_t)$ and for products; computing brackets of Itô processes. \end{quizbox}

\subsection{Consequences of Itô's formula}
\begin{theorem}[Lévy's characterisation] \todo{$\br{X^i,X^j}_t=\delta_{ij}t$ $\Rightarrow$ BM} \end{theorem}
\begin{proofidea} \todo{$\exp(i\xi\cdot X_t+\tfrac12|\xi|^2t)$ is a bounded martingale} \end{proofidea}
\begin{theorem}[Dambis--Dubins--Schwarz] \todo{$M_t=\beta_{\br{M}_t}$} \end{theorem}
\begin{theorem}[Burkholder--Davis--Gundy] \todo{$c_p\E[\br{M}_T^{p/2}]\le\E[(M_T^*)^p]\le C_p\E[\br{M}_T^{p/2}]$} \end{theorem}
\begin{theorem}[Martingale representation] \todo{Brownian filtration: every $L^2$ martingale is $c+\int H\,dB$} \end{theorem}
\begin{usedin} Representation $\Rightarrow$ market completeness / hedging; Lévy $\Rightarrow$ Girsanov proof. \end{usedin}

\subsection{Girsanov's theorem}
\begin{proposition} \todo{$\QQ\ll\PP$ on $\F_\infty$, $D_t=\E[\tfrac{d\QQ}{d\PP}\mid\F_t]$ continuous $\Rightarrow$ $D=\Exp(L)$} \end{proposition}
\begin{theorem}[Girsanov] \todo{$M$ $\PP$-loc.\ mart $\Rightarrow$ $M-\br{M,L}$ is $\QQ$-loc.\ mart; brackets unchanged} \end{theorem}
\begin{proposition}[Novikov] \todo{$\E\exp(\tfrac12\br{L}_\infty)<\infty\Rightarrow\Exp(L)$ UI martingale} \end{proposition}
\begin{quizbox} Girsanov for BM with drift: identify $L$, check Novikov, write the $\QQ$-BM. \end{quizbox}

\subsection{Applications of Girsanov's theorem}
\begin{theorem}[Cameron--Martin] \todo{law of $B+h$, $h\in H$ (Cameron--Martin space)} \end{theorem}
\begin{example} \todo{law of $T_a$ for BM with drift; weak solutions of SDEs with drift} \end{example}
\begin{usedin} Risk-neutral measure; drift removal in pricing. \end{usedin}
